% Diagram pembagian 32 putaran MARS: forward mixing, core, backward mixing.
% Dirujuk dari section-03.tex dengan % Diagram pembagian 32 putaran MARS: forward mixing, core, backward mixing.
% Dirujuk dari section-03.tex dengan % Diagram pembagian 32 putaran MARS: forward mixing, core, backward mixing.
% Dirujuk dari section-03.tex dengan % Diagram pembagian 32 putaran MARS: forward mixing, core, backward mixing.
% Dirujuk dari section-03.tex dengan \input{figures/mars-mixing}.
\begin{figure}[htbp]
    \centering
    \begin{tikzpicture}[
        fase/.style={draw, rounded corners=1pt, minimum height=1.3cm,
                     text width=3.9cm, align=center, font=\small, inner sep=3pt},
        panah/.style={-{Latex[length=2mm]}, thick},
        catatan/.style={font=\footnotesize, align=left, text width=12.9cm},
    ]
        \node[fase] (fwd) at (0,0)
            {\textbf{8 putaran}\\ \textit{forward mixing}\\
             perkalian, rotasi data,\\ penambahan kunci};
        \node[fase] (core) at (4.4,0)
            {\textbf{16 putaran}\\ \textit{core}\\
             kotak-S, rotasi bergantung\\ data, penambahan kunci};
        \node[fase] (bwd) at (8.8,0)
            {\textbf{8 putaran}\\ \textit{backward mixing}\\
             kebalikan \textit{forward mixing},\\ mengacak kembali};

        \draw[panah] (-2.2,0) -- (fwd);
        \draw[panah] (fwd) -- (core);
        \draw[panah] (core) -- (bwd);
        \draw[panah] (bwd) -- (11.0,0);

        \node[catatan, anchor=north] at (4.4,-1.4)
            {Kedua fase \textit{mixing} bekerja tanpa kunci dan dirancang agar tidak
             dapat dibalik. Seluruh 32 putaran memakai kunci hasil perluasan dari
             40 kata, sehingga biaya pemasangan kuncinya paling mahal di antara
             kelima finalis AES};
    \end{tikzpicture}
    \caption{Pembagian 32 putaran MARS. Rancangan ini menempatkan ketahanan terhadap
    kriptanalisis pada putaran inti, sementara putaran tepi mempercepat penyebaran
    perubahan bit masukan}
    \label{fig:mars-mixing}
    \sumber{Munir, \textit{IF4020 Kriptografi}, ITB; Burwick dkk., \textit{MARS}.}
\end{figure}
.
\begin{figure}[htbp]
    \centering
    \begin{tikzpicture}[
        fase/.style={draw, rounded corners=1pt, minimum height=1.3cm,
                     text width=3.9cm, align=center, font=\small, inner sep=3pt},
        panah/.style={-{Latex[length=2mm]}, thick},
        catatan/.style={font=\footnotesize, align=left, text width=12.9cm},
    ]
        \node[fase] (fwd) at (0,0)
            {\textbf{8 putaran}\\ \textit{forward mixing}\\
             perkalian, rotasi data,\\ penambahan kunci};
        \node[fase] (core) at (4.4,0)
            {\textbf{16 putaran}\\ \textit{core}\\
             kotak-S, rotasi bergantung\\ data, penambahan kunci};
        \node[fase] (bwd) at (8.8,0)
            {\textbf{8 putaran}\\ \textit{backward mixing}\\
             kebalikan \textit{forward mixing},\\ mengacak kembali};

        \draw[panah] (-2.2,0) -- (fwd);
        \draw[panah] (fwd) -- (core);
        \draw[panah] (core) -- (bwd);
        \draw[panah] (bwd) -- (11.0,0);

        \node[catatan, anchor=north] at (4.4,-1.4)
            {Kedua fase \textit{mixing} bekerja tanpa kunci dan dirancang agar tidak
             dapat dibalik. Seluruh 32 putaran memakai kunci hasil perluasan dari
             40 kata, sehingga biaya pemasangan kuncinya paling mahal di antara
             kelima finalis AES};
    \end{tikzpicture}
    \caption{Pembagian 32 putaran MARS. Rancangan ini menempatkan ketahanan terhadap
    kriptanalisis pada putaran inti, sementara putaran tepi mempercepat penyebaran
    perubahan bit masukan}
    \label{fig:mars-mixing}
    \sumber{Munir, \textit{IF4020 Kriptografi}, ITB; Burwick dkk., \textit{MARS}.}
\end{figure}
.
\begin{figure}[htbp]
    \centering
    \begin{tikzpicture}[
        fase/.style={draw, rounded corners=1pt, minimum height=1.3cm,
                     text width=3.9cm, align=center, font=\small, inner sep=3pt},
        panah/.style={-{Latex[length=2mm]}, thick},
        catatan/.style={font=\footnotesize, align=left, text width=12.9cm},
    ]
        \node[fase] (fwd) at (0,0)
            {\textbf{8 putaran}\\ \textit{forward mixing}\\
             perkalian, rotasi data,\\ penambahan kunci};
        \node[fase] (core) at (4.4,0)
            {\textbf{16 putaran}\\ \textit{core}\\
             kotak-S, rotasi bergantung\\ data, penambahan kunci};
        \node[fase] (bwd) at (8.8,0)
            {\textbf{8 putaran}\\ \textit{backward mixing}\\
             kebalikan \textit{forward mixing},\\ mengacak kembali};

        \draw[panah] (-2.2,0) -- (fwd);
        \draw[panah] (fwd) -- (core);
        \draw[panah] (core) -- (bwd);
        \draw[panah] (bwd) -- (11.0,0);

        \node[catatan, anchor=north] at (4.4,-1.4)
            {Kedua fase \textit{mixing} bekerja tanpa kunci dan dirancang agar tidak
             dapat dibalik. Seluruh 32 putaran memakai kunci hasil perluasan dari
             40 kata, sehingga biaya pemasangan kuncinya paling mahal di antara
             kelima finalis AES};
    \end{tikzpicture}
    \caption{Pembagian 32 putaran MARS. Rancangan ini menempatkan ketahanan terhadap
    kriptanalisis pada putaran inti, sementara putaran tepi mempercepat penyebaran
    perubahan bit masukan}
    \label{fig:mars-mixing}
    \sumber{Munir, \textit{IF4020 Kriptografi}, ITB; Burwick dkk., \textit{MARS}.}
\end{figure}
.
\begin{figure}[htbp]
    \centering
    \begin{tikzpicture}[
        fase/.style={draw, rounded corners=1pt, minimum height=1.3cm,
                     text width=3.9cm, align=center, font=\small, inner sep=3pt},
        panah/.style={-{Latex[length=2mm]}, thick},
        catatan/.style={font=\footnotesize, align=left, text width=12.9cm},
    ]
        \node[fase] (fwd) at (0,0)
            {\textbf{8 putaran}\\ \textit{forward mixing}\\
             perkalian, rotasi data,\\ penambahan kunci};
        \node[fase] (core) at (4.4,0)
            {\textbf{16 putaran}\\ \textit{core}\\
             kotak-S, rotasi bergantung\\ data, penambahan kunci};
        \node[fase] (bwd) at (8.8,0)
            {\textbf{8 putaran}\\ \textit{backward mixing}\\
             kebalikan \textit{forward mixing},\\ mengacak kembali};

        \draw[panah] (-2.2,0) -- (fwd);
        \draw[panah] (fwd) -- (core);
        \draw[panah] (core) -- (bwd);
        \draw[panah] (bwd) -- (11.0,0);

        \node[catatan, anchor=north] at (4.4,-1.4)
            {Kedua fase \textit{mixing} bekerja tanpa kunci dan dirancang agar tidak
             dapat dibalik. Seluruh 32 putaran memakai kunci hasil perluasan dari
             40 kata, sehingga biaya pemasangan kuncinya paling mahal di antara
             kelima finalis AES};
    \end{tikzpicture}
    \caption{Pembagian 32 putaran MARS. Rancangan ini menempatkan ketahanan terhadap
    kriptanalisis pada putaran inti, sementara putaran tepi mempercepat penyebaran
    perubahan bit masukan}
    \label{fig:mars-mixing}
    \sumber{Munir, \textit{IF4020 Kriptografi}, ITB; Burwick dkk., \textit{MARS}.}
\end{figure}
